\chapter{Introduction}
\label{ch:intro}

\section{Background}

Automated Machine Learning (AutoML) is the subfield of machine
learning that aims to automate the end--to--end construction of
predictive systems, including data cleaning, feature engineering,
model selection, hyper--parameter optimisation and model evaluation.
Early AutoML frameworks such as Auto--Sklearn~\cite{autosklearn} and
AMLB~\cite{amlb} search over a fixed, hand--designed space of
pipelines and return the best configuration found under a budget.
The user describes the task through a structured API or a
configuration file, and the engine returns a trained model and its
measured performance.

The arrival of capable Large Language Models (LLMs) has reshaped this
picture. A user can now describe the task, the data constraints and
the deployment requirements in free--form natural language, and an
LLM--driven agent can decide which features to construct, which
model family to try, how to split the data, and how to evaluate the
result. New systems embed this behaviour into full pipelines,
including feature engineering agents~\cite{caafe},
data--interpreter--style systems~\cite{datainterpreter},
knowledge--driven data--science agents~\cite{dsagent},
tree--search enhanced agents~\cite{sela} and multi--agent
full--pipeline AutoML frameworks~\cite{automlagent}. A recent survey
on LLMs for building and optimising ML workflows~\cite{llmauto}
catalogues this growing landscape and highlights the open problems
that follow from it.

A consequence of this shift is that the automated pipeline now
produces an artifact that traditional AutoML did not surface: a
\emph{plan}, or natural--language justification, that explains what
the system intends to do before it writes a single line of code. An
LLM may, for example, state that it will impute missing values with
the median, use a stratified train/validation split, standardise the
numerical columns using training statistics only, fit a logistic
model under a class--imbalance constraint, and report the macro
F1--score for model selection. The downstream coding agent then
generates the executable pipeline that is supposed to realise this
plan.

Existing checks tell us a great deal about such pipelines: syntax and
type checks reject malformed programs, sandboxed execution surfaces
runtime errors, tests exercise selected behaviour, standard AutoML
evaluation measures predictive quality and cost, and execution
feedback can even be fed back for self--repair~\cite{selfrefine,
critic}. What none of them establishes is whether the program that
runs is the one the preceding plan actually intended.

\section{Problem Statement}

The project targets the following problem.

\begin{quote}
Given an AutoML plan produced by an LLM reasoning agent and a Python
pipeline produced by a coding agent, mechanically determine whether
the pipeline implements that plan, and if it does not, identify
precisely which claims of the plan are missing, contradicted,
mis--parameterised, mis--ordered or merely unverifiable from the
available evidence.
\end{quote}

In the terminology of this project, this is the
\emph{reasoning--code alignment} problem. It is distinct from
reasoning correctness (is the plan itself sound?), from code
correctness (does the code satisfy a code specification?), from
execution correctness (does the code run and produce an output?)
and from predictive performance (does the code reach an acceptable
score?). A plan can be sound and still be realised by an incorrect
program; a program can be well--engineered and still implement a
different plan from the one presented to the user; two programs that
look lexically different may implement the same plan; and a program
that reaches a strong score may still have replaced the intended
estimator or quietly dropped a required pre--processing step.

Formally, if $T$ is the task specification, $R$ the reasoning
artifact, $C$ the corresponding source code, $E$ the execution trace
and $Y$ the achieved performance, then current pipelines check
variants of
\[
Q_R : \text{Is } R \text{ a justifiable strategy for } T?,\quad
Q_C : \text{Does } C \text{ satisfy a code specification?},
\]
\[
Q_E : C \xrightarrow{\text{execute}} E \text{ with satisfactory } Y,
\]
whereas the question of interest here is strictly a fourth one,
\[
Q_A : \text{Does } C \text{ perform the actions specified in } R?
\]
None of $Q_R$, $Q_C$, $Q_E$ implies $Q_A$, which is why
reasoning--code alignment deserves its own verification layer.

\needspace{0.28\textheight}
\section{Motivation}

Three considerations drive the project.

\textbf{Plans are now control artifacts, not just explanations.}
In multi--agent AutoML systems the plan is what a human reviewer or
a downstream module approves; it is also what the user is shown as
the justification for a deployment decision. If the generated code
silently diverges from the plan, every subsequent operational action
is taken on a false premise. A pipeline that is approved on the
basis of a careful, class--imbalance--aware plan but is actually
realised by a random split with a different estimator may still
produce a credible accuracy number, and that number can mask a
systematic failure on the minority class.

\textbf{Execution evidence is weaker than it looks.}
It has already been established that executable code benchmarks
provide a stronger signal than surface similarity for code generation
quality~\cite{humaneval,programsynthesis}, but even passing tests
only verify claimed behaviour and leave latent defects
untouched~\cite{copilot}. A survey of LLMs for code
generation~\cite{jiang2026survey} lists the integration of formal
verification tools within the LLM pipeline as an open safety
direction, rather than a solved problem. Empirical studies of code
generation failure modes~\cite{shen2026failure} characterise which
failures occur most often and how self--critique, multi--model
collaboration or retrieval--augmented generation can enhance
performance; however, all of these frameworks stop at
execution--based failure rather than at reasoning--code alignment.

\textbf{Faithfulness research targets reasoning, not code.}
On the reasoning side, Chain--of--Thought (CoT) prompting is known
to elicit intermediate steps~\cite{wei2023cot}, but these steps do
not have to reflect the model's underlying computation, and recent
work has shown that language models do not always say what they
think~\cite{turpin2023}. Follow--up work has formalised CoT
faithfulness~\cite{young2026lie, jacovi2024chain} and inspected
reasoning chains at a finer granularity through step--level verifier
training and process supervision~\cite{lightman2024}. Typed
Chain--of--Thought~\cite{perrier2025} synthesises reasoning steps
into typed trees that can be validated, and
CRV~\cite{zhao2026crv} builds an attribution graph of model
computation and infers whether the reasoning is correct from
structural fingerprints. These approaches address the correctness of
reasoning rather than whether the emitted AutoML code implements the
reasoning.

The intersection of these three gaps leaves a clear opening: an
AutoML pipeline that passes every available check (compiles, runs,
hits the metric, has a plausible--looking plan) can still be
operationally misleading, and no existing framework is designed to
detect this specific failure mode. This project aims to close that
gap.

\needspace{0.33\textheight}
\section{Objectives}

The objectives of the project are:

\begin{enumerate}
    \item To define reasoning--code alignment as a verification
    question in its own right, kept separate from whether the plan is
    a good strategy, whether the code meets a code specification,
    whether the code runs, and whether it reaches a good score.

    \item To build a working, reproducible pipeline that generates
    AutoML code with an LLM under three prompting conditions (naive,
    schema--aware and meta--feature--guided), runs it in a bounded
    subprocess, and records the plan, the code, the execution outcome
    and the predictive score for every run.

    \item To implement a first, deliberately limited faithfulness
    check that, for each meta--feature--guided run, tests two things
    mechanically: whether every action the plan names actually appears
    in the generated code, and whether every dataset statistic the
    plan cites matches the real meta--features within a small
    tolerance.

    \item To evaluate meta--feature--guided prompting on a controlled
    270--cell sweep (6 datasets $\times$ 3 conditions $\times$ 3 LLM
    backends $\times$ 5 seeds), measuring runnability, iteration count,
    error mix, predictive score and faithfulness, and to run the
    faithfulness check on the real meta--feature runs.

    \item To state honestly what this first check cannot yet do, and
    to describe --- as a proposed design and future work, not a
    delivered system --- the fuller claim--by--claim verifier that
    would be needed to close the gap.
\end{enumerate}

\section{Scope}

The project targets \emph{compositional Python tabular classification
and regression} pipelines built with scikit--learn~\cite{sklearn}
and its \texttt{Pipeline}, \texttt{ColumnTransformer}, estimator,
scoring and custom--transformer abstractions. Within this setting,
the \emph{implemented} faithfulness check tests two things for each
run: whether each operator or estimator named in the plan (for
example a scaler, an imputer or a specific classifier) appears in the
generated code, and whether each dataset statistic the plan cites
(for example a missing--value ratio or a class--balance figure)
matches the value actually computed from the data. A \emph{fuller}
verifier --- designed in Chapter~\ref*{ch:method} but not implemented
in this project --- would additionally check:
\begin{itemize}
    \item the columns to which each operator is applied,
    \item the parameters and ordering of the pipeline,
    \item leakage--relevant properties such as whether a scaler is
    fitted before or after the train/validation split,
    \item evaluation metrics and class--balance constraints.
\end{itemize}
These richer checks are stated as a target, not as a delivered
capability.

The following items are explicitly \emph{out of scope} for the first
version of the system:
\begin{itemize}
    \item Deep--learning pipelines and distributed training systems;
    multi--modal workflows; notebook pipelines with hidden state;
    AutoML frameworks other than scikit--learn--style tabular
    pipelines.
    \item Verification of the plan itself (whether the chosen
    strategy is a good strategy for the task) --- the verifier only
    checks implementation faithfulness relative to a given plan, not
    the plan's merit.
    \item Mechanistic faithfulness in the sense of
    CRV~\cite{zhao2026crv}: the project checks \emph{artifact--to
    --artifact} conformance (plan $\to$ emitted code), not conformance
    between the model's internal computation and its emitted
    reasoning.
    \item Automatic repair of misaligned pipelines; the verifier is
    designed to report, not to patch.
\end{itemize}

\section{Contributions}

The contributions of the project are of three kinds --- conceptual,
engineering and empirical.

\begin{enumerate}
    \item \textbf{(Conceptual.)} A clear separation of reasoning--code
    alignment from reasoning correctness, program correctness,
    execution success and predictive performance, stated as the
    fourth question $Q_A$ above.

    \item \textbf{(Conceptual, proposed.)} A design
    (Chapter~\ref*{ch:method}) for a fuller verifier: a structured
    form of the plan and of the code so that each plan claim can be
    matched to the part of the code that realises it, together with a
    mismatch taxonomy for AutoML pipelines. This design is specified,
    not implemented, and is not evaluated in this report.

    \item \textbf{(Engineering.)} A running, reproducible pipeline ---
    five FastAPI backend services (data, meta--features, generation,
    analysis and gateway) and a Streamlit front--end, a Redis--backed
    job queue and worker, a
    configurable LLM runner (Ollama cloud models in the reported
    runs), a bounded subprocess executor and a PostgreSQL database with
    local--filesystem artefact storage --- that records the plan, code,
    execution trace and score of every run.

    \item \textbf{(Engineering.)} A first mechanical faithfulness
    check on the project's \texttt{main} branch that, for every
    meta--feature--guided run, tests each cited action against the
    generated code and each cited statistic against the real
    meta--features, and returns a single pass/fail verdict for the run.

    \item \textbf{(Empirical.)} A controlled 270--cell study of
    meta--feature--guided prompting, measured along five axes
    (runnability, iteration cost, error mix, predictive score and
    faithfulness), providing evidence towards the research questions
    RQ1--RQ5 (defined in Chapter~\ref*{ch:lit}). Running the check on
    the real meta--feature runs flagged seven pipelines that executed
    and scored well but whose stated reasoning did not match their
    code.
\end{enumerate}

\paragraph{Implementation and validation status.}
To keep the claims of this report honest, the work is separated into
three levels throughout. \emph{Designed} items --- the fuller
claim--by--claim verifier, the mismatch taxonomy and the baseline
comparison protocol --- are specified in Chapters~\ref*{ch:method}
and~\ref*{ch:design} but are not implemented in this project.
\emph{Implemented} items --- the five backend services and the
Streamlit front--end, the three prompt conditions, the bounded
subprocess executor and the first faithfulness check --- exist and run
on the
project's
\texttt{main} branch. \emph{Validated} items --- the prompting study
and the first faithfulness check --- are supported by the 270--cell
sweep in Chapter~\ref*{ch:results}. The report never presents a
designed item as though it were implemented or validated, and every
number in Chapter~\ref*{ch:results} is traceable to the recorded
experimental run.

\section{Organization of the Report}

The remainder of this report is organised as follows.

\textbf{Chapter~\ref*{ch:lit}, Literature Review}, surveys
LLM--based AutoML, code and reasoning verification, formal typed
and graph--based verification, and identifies the research gap that
motivates the project.

\textbf{Chapter~\ref*{ch:method}, Proposed Methodology}, defines
the four correctness questions and the reasoning--code alignment
relation, and sets out the \emph{proposed} design for a fuller
verifier --- structured forms of the plan and the code, the matching
rules, the mismatch taxonomy and the associated measures --- which is
specified but not implemented in this project.

\textbf{Chapter~\ref*{ch:design}, System Design}, presents the
end--to--end architecture, the execution and evaluation pipeline,
the module decomposition, the data--flow and the storage schema.

\textbf{Chapter~\ref*{ch:impl}, Implementation}, describes the
development environment, the implementation of each module, the
integration with the Ollama LLM runner, the bounded subprocess
execution layer and the storage of faithfulness verdicts.

\textbf{Chapter~\ref*{ch:results}, Results and Discussion},
reports the experimental protocol and the 270--cell sweep, its
results across the five research questions (runnability, iteration
cost, error mix, predictive score and faithfulness), the statistical
analysis, and a discussion of what the faithfulness check did and did
not catch.

\textbf{Chapter~\ref*{ch:conclusion}, Conclusion and Future
Scope}, summarises the work, states the limitations and threats to
validity, and outlines realistic extensions of the framework.

The appendix presents supporting material and the mapping of the
project to Course Outcomes, Program Outcomes and Programme Specific
Outcomes.
